\documentclass[10pt]{amsart}
\usepackage[a4paper,margin=26mm]{geometry}
\usepackage{amssymb,mathtools,hyperref,microtype}
\hypersetup{hidelinks}\emergencystretch=1em
\newcommand{\one}{\mathbf1}\newcommand{\C}{\mathcal C}
\newcommand{\D}{\mathcal D}\newcommand{\E}{\mathcal E}
\newcommand{\Hom}{\operatorname{Hom}}\newcommand{\Irr}{\operatorname{Irr}}
\newcommand{\tr}{\operatorname{tr}}\newcommand{\ptr}{\operatorname{ptr}}
\newcommand{\im}{\operatorname{im}}
\newtheorem{theorem}{Theorem}[section]
\newtheorem{lemma}[theorem]{Lemma}
\newtheorem{corollary}[theorem]{Corollary}
\newtheorem{question}[theorem]{Question}
\title{Joint monodromy and finite gauging}
\author{Tethys}
\author{Capella}
\date{22 November 2026}
\begin{document}\raggedbottom
\begin{abstract}
We prove that finite joint monodromy detection is preserved by
faithful finite gauging and by Deligne products of braided fusion
categories. The closure theorem requires neither unitarity nor a
ribbon structure. Categorical partial traces provide quantitative
criteria for the required fixed subobjects to vanish in unitary
modular categories. These criteria apply to every unitary
categorification of the specified extended-Haagerup and level-four
$F_4$ modular data. An exclusive-vacuum reconstruction theorem
then gives all-place braid lattices, integral framed link values,
and good modular ribbon reduction away from the twist order for
their iterated products and finite gaugings, and their coefficient
conjugates. Prime permutation orbifolds can eliminate every
contracting simple probe even when joint detection persists.
\end{abstract}\maketitle

\section{Introduction}
The monodromy of a distinguished object with a single probe
can have fixed vectors even when that object is not transparent.
For a tensor word of probes the relevant condition is joint:
the distinguished point must be fixed under every point push
around an input strand. We use the common fixed \emph{subobject},
so that vanishing persists after arbitrary spectator objects
are inserted. This supplies a condition compatible with finite
gauging, where the spectators may belong to nonneutral sectors.

Theorem \ref{closure} constructs a detecting word after an
arbitrary faithfully graded finite crossed extension and
equivariantization. Electric representations exclude distinguished
nonidentity degrees. Fluxes then transport the neutral probe
components independently, reducing their joint constraints to
those in the neutral category. The final electric outputs are
excluded by generation of the finite group. These steps require
invertible transport and semisimplicity; they do not use an
inner product or positivity.

The quantitative results use positivity at a different point.
The trace of a palindromic product of monodromy projections
bounds the dimension of their common image. Repetition gives
vanishing when one matched-channel dimension is small enough.
A dimension gap can already force a two-probe common image
into copies of the unit, where different probe twists or an
additional meridian eliminate it. None of these estimates
requires multiplicity-free fusion.

The arithmetic consequence uses the exclusive-vacuum theorem
\cite[Theorem 17]{P56}: one full-support coloured word with a
vacuum constituent occurring at no other simple output implies
the integral braid and ribbon statements. The present construction
provides that hypothesis uniformly across categorifications of
specified modular data, without their individual $F$- and
$R$-symbols. The extended-Haagerup example and its finite
gaugings are not subject to a finite braid-image hypothesis.
Finite images for weakly group-theoretical braided categories
are already known \cite[Theorem 7.2]{GN}; no new finiteness
claim for that class is intended here.

Theorem \ref{permutation-main} shows that the one-probe
dimension criterion does not survive gauging. Its hypothesis
requires only a noninvertible boson with a nonempty monodromy-one
summand against every simple probe in a completely rational net
category. Positive asymptotic tensor return follows from that
hypothesis by modular diagonalization. No self-duality or separate
tensor-return hypothesis is needed.

All categories are complex linear and fusion categories are split.
When a ribbon structure is specified, it is part of the data.
The results have been checked by AI researchers and have not been
checked by a human mathematician.

\section{Point pushes and finite gauging}
Let $\C$ be a complex braided fusion category, let $a$ be simple,
and let $v=(x_1,\ldots,x_n)$ be a word with tensor product $V$.
In the braid group with one distinguished strand of colour $a^*$
followed by $v$, forgetting the distinguished strand has kernel
\[
             F_n=\pi_1(\mathbb C\setminus\{n\text{ points}\}),
\]
generated by moving that strand positively around each input
along a fixed star system of paths. This is the standard
forgetting-strand exact sequence; see
\cite[Theorem 1.16 and Section 1.5.2]{KT}.
Define $I_a(v)\subset a^*V$ to be its common fixed subobject:
on $\Hom(t,a^*V)$ it is the common fixed subspace for these
point pushes, for every simple $t$.
Semisimplicity gives this subobject by intersecting the kernels
of the generators minus identity.

\begin{lemma}\label{normal}
The subobject $I_a(v)$ is invariant under the colour-preserving
input braid group. Input braid isomorphisms $V\to V'$ carry it
to $I_a(v')$. If $w$ contains $v$ as a designated subword and
$I_a(v)=0$, then $I_a(w)=0$.
\end{lemma}
\begin{proof}
The subgroup $F_n$ is normal in the group of braids fixing
the distinguished strand. Its fixed subobject is therefore
preserved by the input braid section. The same normality gives
the groupoid statement for reorderings of the input colours.
After such a reordering put the designated subword first.
Its point pushes act on $a^*V$ tensored with the remaining
word; they are among the point pushes for the larger word.
Their fixed subobject is $I_a(v)$ tensored with that word.
The larger fixed subobject is contained in it.
\end{proof}

A word annihilating $I_a$ for every nonunit simple exists
if one has a separate word for each such simple: concatenate
them and apply Lemma \ref{normal}. The condition is stronger
than purity of one sphere quotient, since it concerns an
object, rather than just its maps from the unit.

\begin{theorem}\label{closure}
Let $\D$ be a complex braided fusion category. Suppose a word $v$ of
its simple objects satisfies
\[
                         I_a(v)=0\quad(a\in\Irr(\D),\ a\ne\one).
                                                               \tag{G1}
\]
Let $\mathcal T=\bigoplus_{g\in G}\mathcal T_g$ be any faithfully
graded complex $G$-crossed braided extension with
$\mathcal T_e=\D$, for a finite group $G$.
Then its equivariantization $\E=\mathcal T^G$ has a word $w$
with $I_A(w)=0$ for every nonunit simple $A$ of $\E$.
Condition (G1) is also preserved by finite Deligne products.
\end{theorem}

We construct $w$ explicitly. Include a family of irreducible
electric objects $V_\pi\in\operatorname{Rep}G\subset\E$
whose kernels intersect trivially. Include one simple flux
object supported on each conjugacy class of $G$.
Finally, for each occurrence of $x$ in $v$, include one simple
neutral object of $\E$ whose underlying object contains $x$.
Such an object exists by decomposing induction
$\bigoplus_{g\in G}g(x)$.
Its underlying simple labels are the whole $G$-orbit of $x$.
A flux object exists for every class because the grading is faithful.
Order the designated neutral probes before the flux probes.

We use that any selection of representatives of all the
conjugacy classes generates $G$. For if a proper subgroup met
every class, its transitive action on its cosets would have a
fixed point for every group element. The average number of
fixed cosets is one, whereas the identity has more than one
and all other elements would have at least one. 

\subsection{Two crossed-braid calculations}
Write $u_g^X:g(X)\to X$ for the equivariant structure of $X$.
For homogeneous $Y\in\mathcal T_g$ the braiding in $\E$ is
\[
 c^\E_{Y,X}=(u_g^X\otimes1_Y)c^\times_{Y,X},
 \qquad c^\times_{Y,X}:YX\longrightarrow g(X)Y.          \tag{G2}
\]
We suppress the specified associators and action comparisons.
Units may be normalized so that the action on $\one$ and
crossed braiding with $\one$ are the unit maps.
The electric object $V_\pi$ has underlying object
$\one^{\oplus\dim\pi}$, with $u_g^{V_\pi}=\pi(g)$.
These are the usual equivariantization conventions of
\cite[Sections 3.2--3.4]{CGPW}.

Decompose the underlying objects of the simple $\E$-factors
into their simple $\mathcal T$-isotypes, retaining their finite
multiplicity spaces. Their tensor product is an external
direct sum indexed by the tuple of chosen component labels.
Call one such summand a component summand. These summands may
be isomorphic; their external labels are nevertheless fixed.
We do not assume their projections belong to a braid-image algebra.

\begin{lemma}\label{electric}
Let an electric input $V_\pi$ make a pure loop around input $j$,
along a path avoiding the other inputs, and return to its position.
On a component summand it acts only on $V_\pi$, by $\pi(s_j)$,
where $s_j$ is conjugate to the chosen degree of input $j$.
It preserves every component label and every other internal vector.
If $V_\pi$ is adjacent to a distinguished homogeneous object
of degree $g$, its point-push monodromy is $\pi(g)$ on $V_\pi$.
\end{lemma}
\begin{proof}
For adjacent factors, (G2) gives one identity crossing on the
electric multiplicity space and one action of the flux degree.
The crossed braidings with its underlying unit copies are the
unit maps, so no operator remains on any other factor.
For a nonadjacent path the crossings along the path and its
reverse contribute inverse matrices on the same electric factor;
they conjugate $\pi(g_j)$ by a product of crossed degrees.
All other factor maps cancel. This proves the first assertion.
The same adjacent calculation proves the second.
This argument uses no three-cocycle triviality.
\end{proof}

\begin{lemma}\label{neutral}
Let a neutral input $U\in\E$, with a chosen underlying simple
component $x\in\D$, make the same pure loop around input $j$.
The braid permutes only the component label of $U$, taking
$x$ to $s_j(x)$, where $s_j$ is conjugate to the chosen degree
of input $j$. It fixes all other component labels and all
degrees. On each affected summand it is a isomorphism.
The elements $s_j$ depend only on the chosen degree tuple and
the paths, not on the neutral component label.
\end{lemma}
\begin{proof}
For an adjacent homogeneous factor $Y_h$, the first crossing
has component labels $(x,Y_h)\mapsto(Y_h,x)$, since $x$ has
degree $e$. The second has labels
$(Y_h,x)\mapsto(h(x),Y_h)$ by (G2).
Thus only $x$ changes. Both arrows are isomorphisms, including
the equivariance map on its multiplicity space.

Along a path past other inputs the moving factor always
has degree $e$. Their degrees and component labels therefore
remain fixed under the crossed action. The path and reverse
path conjugate $h$ by the product of crossed degrees, giving
$s_j$. They need not cancel on internal fusion vectors,
and no such cancellation is asserted. Their composition is
still a isomorphism between exactly the displayed
component summands. The label calculation depends only on
degrees, which proves the last assertion.
\end{proof}

\subsection{Proof of the gauging theorem}
We apply the component calculations of the preceding section.
Fix a nonunit simple $A$ of $\E$ and write $K=I_A(w)$.
It is invariant under input braids by Lemma \ref{normal}.
All arguments can be tested on $\Hom(t,UA^*\,UW)$ for each
simple $t$ of $\mathcal T$, where $U$ is the forgetful functor;
thus they concern the whole object $UK$.

First suppose a component of $UA^*$ has degree $g\ne e$.
Choose an electric probe $\pi$ with $\pi(g)\ne1$, and
reorder it next to $A^*$ using Lemma \ref{normal}.
The adjacent distinguished point push fixes $K$ and acts
by $\pi(g)$ on that electric factor.

The input loops of this electric factor around the selected
flux factors preserve all component summands and act by
$\pi(s_j)$, by Lemma \ref{electric}. The $s_j$ still represent
all conjugacy classes, so they generate $G$.
On a fixed component summand the projection of $K$ is
therefore stable under $\pi(G)$ on this factor and is
annihilated by $\pi(g)-1$.
Irreducibility of $\pi$ makes this space zero: a
$\pi(G)$-stable subspace of $V_\pi\otimes L$ is
$V_\pi\otimes L'$; it cannot be annihilated by the nonzero
operator $\pi(g)-1$ unless $L'=0$.
The component projection is legitimate here because every
operator just used preserves every component summand.
This excludes every distinguished nonidentity degree.

Now suppose a nonzero component of $UA^*$ is $a^*\in\D$,
with $a\ne\one$. For each designated neutral probe lift
corresponding to $x_i$ in $v$, its chosen component is some
translate of $x_i$. The input loops in Lemma \ref{neutral}
act on this one label through elements $s_j$ generating $G$.
They fix all other component labels. Thus a sequence of such
input braids can independently bring every designated
component to exactly $x_i$.
The component label of $A^*$ is unchanged because these
braids act only on the inputs.

A braid gives a bijection on the tuple labels and an
isomorphism on each summand. Hence a nonzero component of
a vector from $UK$ remains nonzero after this transport;
different summands cannot cancel there.
Reorder the designated neutral probes first.
When moving a neutral probe $U$ left past $Y$, choose
$(c_{U,Y})^{-1}$. Its degree is $e$, so this crossing
preserves both chosen component labels. Equivalently one
can put the designated probes first before their transport.
On the transported component, their distinguished point
pushes are exactly the base point pushes on $a^*V$,
tensored with identity on their multiplicity spaces and
on the remaining word, even if that word has flux factors.
Their common fixed subobject is zero by (G1).
This contradicts nonzero transport and excludes every
distinguished nonunit neutral component.

A remaining simple $A$ has underlying object a sum of units.
It is therefore an electric irreducible $V_\rho$, with
$\rho$ nontrivial because $A\ne\one$.
The distinguished point pushes around the selected flux
inputs act only by $\rho^*(s_j)$ on this distinguished
multiplicity factor, as in Lemma \ref{electric}.
Their common fixed space is
$(V_{\rho^*})^{\langle s_j\rangle}=(V_{\rho^*})^G=0$.
This excludes the last possibility, so $K=0$.

For Deligne products, concatenate the base detecting words
embedded one factor at a time. If a simple tuple is nonunit,
some factor $a_i$ is nonunit, and its embedded probe word
has fixed subobject $I_{a_i}(v_i)$ tensored with the other
factors of the distinguished tuple. It is zero. Lemma
\ref{normal} proves (G1) for the whole concatenation.
This completes the proof of Theorem \ref{closure}.


\section{Vacuum constituents and arithmetic}
Let $\C$ be a complex modular ribbon category. For a word
$w=(x_1,\ldots,x_n)$ with product $W$, let $A_w$ be the image
of its colour-preserving braid group in $\operatorname{End}(W)$.
Let $\Omega$ be the positive full twist and let $M_i$ be the
pure loop moving input $i$ positively around all other inputs,
with no individual framing twists. Define
\[
 J_w=\left(\Omega-\Bigl(\prod_i\theta_{x_i}\Bigr)^{-1},\quad
                    M_i-\theta_{x_i}^{-2}\ (1\le i\le n)\right)
                    \triangleleft A_w,
 \qquad Q_w=W/J_wW.
\]
The ideal is two-sided. The braid-image algebras are semisimple
by \cite[Theorem 1.1]{GS}; this includes the colour-preserving
group, a finite-index extension of the pure group.
Thus a quotient multiplicity space is realized by a
subrepresentation annihilated by the ideal.

\begin{lemma}\label{sphere-bending}
For a simple output $a$, the full-twist relation annihilates
its whole multiplicity unless $\theta_a=1$. In that case,
rigid bending identifies a surviving multiplicity with vectors
in $\Hom(\one,a^*W)$ fixed by the distinguished point pushes.
The whole vacuum multiplicity survives the ideal.
\end{lemma}
\begin{proof}
Balancing makes the full twist on this output the scalar
$\theta_a(\prod_i\theta_{x_i})^{-1}$, as in book tag
\texttt{Z7OZ}. Put one chosen input $x$ first and write
the remaining word as $Y$. On a simple $b$-channel of
$\Hom(a,xY)$, its sphere loop has eigenvalue
$\theta_a/(\theta_x\theta_b)$. Bending pairs the
$a^*x$ channel with $Y$, and its monodromy has eigenvalue
$\theta_b/(\theta_a\theta_x)$. Consequently that
distinguished double crossing is $\theta_x^{-2}M_i^{-1}$
under the rigid identification. The sphere relation makes
it identity. The same ribbon isotopy identifies the operators,
including their multiplicity spaces. The two-sided ideal
imposes the conjugate peripheral paths, and reorderings
transport the ideal to the corresponding word. Products
and inverses are fixed as well. The distinguished framing
adds no scalar because $\theta_{a^*}=\theta_a=1$.
At the unit output the balancing relations are precisely
the displayed scalar relations, so every vacuum vector survives.
\end{proof}

\begin{corollary}\label{arithmetic}
Let $\C$ be a complex modular ribbon category of twist order
$N$. Suppose a finite word $v$ satisfies $I_a(v)=0$ for every
nonunit simple $a$. Every simple-coloured braid representation
of $\C$ has a full invariant lattice at every finite place
after finite local coefficient extension. Every simple-coloured
framed link value is an algebraic integer. At each finite
place of residue characteristic prime to $N$, $\C$ has a
split integral ribbon form with modular special fibre and
the same labelled fusion rules after finite local extension.
These assertions also hold for every coefficient conjugate.
\end{corollary}
\begin{proof}
Enlarge $v$ by its dual word and by a simple/dual pair for
every simple colour. The product has a nonzero unit summand,
and Lemma \ref{normal} preserves every zero fixed-subobject
test. Lemma \ref{sphere-bending} shows that the sphere
quotient is entirely vacuum and retains the whole vacuum.
Any irreducible constituent of this vacuum is exclusive
to the unit output: if it occurred at another output,
the same two-sided ideal would annihilate it there.
The enlarged word has every colour, so
\cite[Theorem 17]{P56} gives all three conclusions.
Exclusivity is an algebraic representation property and
is preserved by coefficient embeddings, proving the last
assertion without unitarity of the conjugates.
\end{proof}

Theorem \ref{closure} and Corollary \ref{arithmetic} apply
after any finite sequence of faithful finite gaugings and
Deligne products for which the resulting category has a
specified modular ribbon structure. The braided closure
theorem does not choose that structure.
Taking the neutral category to be $\mathrm{Vec}$ recovers
all twisted finite-group doubles, since its detecting
condition is empty and a faithful crossed extension has
underlying pointed category $\mathrm{Vec}_G^\omega$.


\section{Quantitative joint detection}
\subsection{A partial-trace identity}
Traces in this section are categorical spherical traces. Let $c$ denote
the braiding and suppress associators. The following is the
partial-trace superposing identity, with its conventions spelled out.

\begin{lemma}\label{trace-superpose}
Let $A,X,Y$ be objects of a complex ribbon category and let
$f\in\operatorname{End}(A Y)$. Put
\[
 C=1_A\otimes c_{X,Y}:A X Y\longrightarrow A Y X,
 \qquad \widehat f=C^{-1}(f\otimes1_X)C.
\]
Then $\ptr_Y(\widehat f)=\ptr_Y(f)\otimes1_X$.
\end{lemma}
\begin{proof}
Write $b_Y:\one\to YY^*$ for the left coevaluation and
$\widetilde e_Y:YY^*\to\one$ for the right evaluation.
Naturality and the hexagon at these two morphisms give
\begin{align*}
 (c_{X,Y}\otimes1_{Y^*})(1_X\otimes b_Y)
   &=(1_Y\otimes c_{X,Y^*}^{-1})(b_Y\otimes1_X),\\
 (1_X\otimes\widetilde e_Y)(c_{X,Y}^{-1}\otimes1_{Y^*})
   &=(\widetilde e_Y\otimes1_X)(1_Y\otimes c_{X,Y^*}).
\end{align*}
Substitute them into
$\ptr_Y(\widehat f)=(1_{AX}\otimes\widetilde e_Y)
(\widehat f\otimes1_{Y^*})(1_{AX}\otimes b_Y)$.
The two resulting crossings on $X Y^*$ commute past
$f$ on $A Y$ and cancel. The remaining expression is
$\ptr_Y(f)\otimes1_X$. The argument retains the specified
pivotal evaluation, so it is a categorical partial trace.
\end{proof}

Let $A$ be simple of dimension $d_A\ne0$, and choose idempotents
$p_j\in\operatorname{End}(A X_j)$.
On $T_m=A X_1\cdots X_m$, transport $p_j$ to $A X_j$
by crossing $X_j$ positively past $X_{j-1},\ldots,X_1$,
apply $p_j$, and undo that path. Denote the resulting
endomorphism by $P_j$, extending it by identities on later factors.
It acts only on the first $j+1$ factors.
For $j=1$, this is $p_1$ on the first two factors.
Let $D_j=\tr_{A X_j}(p_j)$ and $\lambda_j=D_j/d_A$.

\begin{lemma}\label{trace-palindrome}
For
\[
 F_m=P_1P_2\cdots P_{m-1}P_mP_{m-1}\cdots P_2P_1
\]
one has $\tr_{T_m}(F_m)=d_A\prod_{j=1}^m\lambda_j$.
\end{lemma}
\begin{proof}
The partial trace of $p_m$ is scalar on simple $A$. Its full trace
is $D_m$, so the scalar is $\lambda_m$. Repeated application of
Lemma \ref{trace-superpose}, with spectator $X_1\cdots X_{m-1}$,
gives $\ptr_{X_m}(P_m)=\lambda_m1_{T_{m-1}}$.
Every other factor in $F_m$ acts on $T_{m-1}$, so naturality of
partial trace and $P_{m-1}^2=P_{m-1}$ give
$\ptr_{X_m}(F_m)=\lambda_mF_{m-1}$.
For $m=1$ the trace is $D_1=d_A\lambda_1$.
Induction proves the formula.
\end{proof}

\subsection{The unitary estimate}
Now let $\C$ be a unitary ribbon fusion category with its positive
spherical structure. Its simple dimensions $d_z$ are positive and
satisfy $d_z\ge1$, since $z z^*$ contains the unit.
Assume $p_j$ are orthogonal projections. Then their braid transports
$P_j$ are orthogonal projections. Write
$I_m=\bigcap_{j=1}^m\im P_j$, an intersection of subobjects of $T_m$.

\begin{theorem}\label{trace-bound}
One has
\[
                    d_{I_m}\le d_A\prod_{j=1}^m\lambda_j.    \tag{T1}
\]
In particular $I_m=0$ if the right side is smaller than one.
\end{theorem}
\begin{proof}
The endomorphism $F_m$ is positive: it is $B^*P_mB$ for
$B=P_{m-1}\cdots P_1$. The common image $I_m$ and its orthogonal
complement are preserved by every $P_j$. On $I_m$, $F_m$ is identity.
On each simple isotype $z^{\oplus r}$, categorical trace is $d_z$
times ordinary matrix trace. Thus positivity on the complement
gives $\tr(F_m)\ge d_{I_m}$. Lemma \ref{trace-palindrome} proves (T1).
Every nonzero object has dimension at least one, giving the last claim.
\end{proof}


For a simple $a$ in a unitary modular category, let
\[
 D_a(x)=\sum_{c:\,\theta_c=\theta_a\theta_x}N_{a^*,x}^c d_c
\]
be the dimension of its monodromy-one summand against $x$.
For a boson the condition reduces to $\theta_c=\theta_x$.
The spectral projectors are orthogonal, with no restriction
on the fusion multiplicities or the number of channels.

\begin{corollary}\label{repeat-detect}
If $D_a(x)<d_a$, then $I_a(x^m)=0$ for every $m$ satisfying
$d_a(D_a(x)/d_a)^m<1$.
\end{corollary}
\begin{proof}
Apply Theorem \ref{trace-bound} to the monodromy-one
projectors. The common fixed subobject is their common image.
\end{proof}


\subsection{Dimension gaps and meridian compression}
In this section let $\C$ be unitary modular, let $H$ be its positive Hopf matrix with $H_{0,0}=1$, and fix a simple boson $a$. Fix two probes $x,z$, set $A=a^*$ and $T=Axz$, and put
\[
 \delta(T)=\min\{d_b:b\ne\one,\ \Hom(b,T)\ne0\},
\]
with minimum $+\infty$ if this set is empty.

\begin{lemma}\label{dimension-gap}
Let $P_1,P_2$ be the monodromy-one projectors on $Ax$ and
on $Az$, the second transported by crossing $z$ positively
past $x$ and undoing that path. If
\[
                  \frac{D_a(x)D_a(z)}{d_a}<\delta(T),            \tag{D1}
\]
then $I=\im P_1\cap\im P_2$ is a sum of copies of the unit.
If $\Hom(\one,T)=0$ or $\theta_x\ne\theta_z$, then $I=0$.
\end{lemma}
\begin{proof}
These are orthogonal projections. The partial-trace superposing
identity of Lemma \ref{trace-superpose}, gives
$\ptr_z(P_2)=(D_a(z)/d_a)1_{Ax}$, and therefore
\[
 \tr_T(P_1P_2P_1)
 =\frac{D_a(z)}{d_a}\tr_{Ax}(P_1)
 =\frac{D_a(x)D_a(z)}{d_a}.
\]
The positive operator $P_1P_2P_1$ is identity on $I$.
Its trace on the orthogonal complement is nonnegative,
so $d_I$ is at most the displayed trace.
A nonunit simple summand of $I$ would contribute at least
$\delta(T)$, contradicting (D1). If $T$ has no unit summand,
then $I=0$. On its unit isotype the intermediate object in $Ax$
is $z^*$, so balancing makes $M_{A,x}$ the scalar
$\theta_z/\theta_x$, including every multiplicity.
If the two twists differ, $P_1$ annihilates that whole isotype,
and again $I=0$.
\end{proof}

\subsection{Compression on an arbitrary unit multiplicity}
Let $y$ be another simple probe and let $E\in\operatorname{End}(T)$
be the orthogonal projector onto the whole unit isotype of $T$.
Write $ET\simeq\one^{\oplus k}$. Transport the monodromy $M_{A,y}$
past the spectators $xz$:
\[
 C=1_A\otimes c_{xz,y}:Axzy\longrightarrow Ayxz,\qquad
 B=C^{-1}(M_{A,y}\otimes1_{xz})C.
\]
\begin{lemma}\label{meridian-compression}
On $ET\otimes y$, one has
\[
           (E\otimes1_y)B(E\otimes1_y)
       =\frac{H_{A,y}}{d_a d_y}\,1_{ET\otimes y}.                \tag{D2}
\]
This holds for every $k$, including zero.
\end{lemma}
\begin{proof}
Lemma \ref{trace-superpose}, applied to the coupon $M_{A,y}$ and
spectator $xz$, gives
$\ptr_y B=(H_{A,y}/d_a)1_T$; the scalar is identified by
taking the full Hopf trace.
Compression to $ET\simeq\one^{\oplus k}$ makes the left side
of (D2) a matrix $K$ on its multiplicity space tensored with
$1_y$, because $y$ is simple. Its partial trace is $d_yK$.
Naturality of partial trace also identifies it with
$(H_{A,y}/d_a)1_{ET}$. This proves (D2).
\end{proof}


The dimension-gap argument also gives an object-level test.
With the notation of Lemma \ref{dimension-gap}, different probe
twists make $I_a(x,z)=0$. If their twists agree, adjoining
a probe $y$ with $H_{a^*,y}\ne d_a d_y$ gives
$I_a(x,z,y)=0$: the first two constraints put the common
image inside the full unit isotype of $a^*xz$; compressing
the third fixed-vector equation by its projector makes
it $v=(H_{a^*,y}/d_a d_y)v$ on every simple multiplicity
space, by Lemma \ref{meridian-compression}. This proves
vanishing as an object and hence after arbitrary spectators.


\section{Numerical criteria and the two nonintegral seeds}\label{sec:seeds}
Use the labelled full positive records B22-00332 and
B16-00499 of the Anyon Atlas \cite{Atlas}. Exact fusion
coefficients and twists refer to its supplied order.
The following are finite inputs, reconstructed from its
cyclotomic Hopf matrices by the accompanying programs.
\[
\begin{array}{c|c|c|c}
\text{record}&\text{simple twist orders}&\text{nonunit bosons}
                           &\text{probes of coprime orders}\\\hline
\text{B22-00332}&1,5,13&5,6,7,8,21&\operatorname{ord}(\theta_1)=13,
                                      \operatorname{ord}(\theta_9)=5\\
\text{B16-00499}&1,2,13&5,10&\operatorname{ord}(\theta_1)=13,
                                      \operatorname{ord}(\theta_{15})=2
\end{array}
\]
Neither record has a simple twist of the product order
$65$ or $26$, respectively. Thus every nonbosonic $a$
has $D_a(x)=0$ for the probe of the other nontrivial
order: balancing would require a twist of that product order.

For B22-00332 the nonunit dimensions are greater than $48$,
and
\[
 48<d_1,d_2<97/2,\quad 49<d_5<50,\quad
 114<d_j<115\ (j=6,7,8),
\]
\[
 D_{X_j}(X_1)=d_1,\quad D_{X_j}(X_2)=d_2\quad(j=5,6,7,8),
 \qquad D_{X_{21}}(X_1)=0.
\]
For $j=5,6,7,8$ one has $d_1d_2/d_j<48$.
Moreover $\theta_1=\exp(2\pi i/13)$ and
$\theta_2=\exp(10\pi i/13)$ are distinct. Lemma
\ref{dimension-gap} therefore gives $I_{X_j}(X_1,X_2)=0$
for every nonunit boson. Combined with the nonboson
tests, the three-colour word $(X_1,X_2,X_9)$ satisfies
the hypothesis of Theorem \ref{closure} for every unitary
categorification of the specified full datum.
The same two-probe test shows that the sphere quotient of
$(X_1,X_2,X_3,X_4)$ is $\one^{\oplus17}$, since
$\sum_bN_{1,2}^bN_{3,4}^{b^*}=17$.

For B16-00499, $D_{X_5}(X_1)=d_1$ and
$D_{X_{10}}(X_1)=0$, with $d_1<8$ and $18<d_5<19$.
Since $19(8/18)^4<1$, Corollary \ref{repeat-detect}
annihilates both nonunit bosons with four copies of $X_1$.
The word $(X_1^4,X_{15})$ consequently satisfies the same
all-output hypothesis in every unitary categorification.

\begin{corollary}\label{seed-arithmetic}
Every unitary categorification of either specified full
positive datum, every category obtained from them by
finite Deligne products and faithful finite unitary gaugings,
and every coefficient conjugate, has all the arithmetic
properties of Corollary \ref{arithmetic}, with the twist
order of the resulting specified ribbon category.
\end{corollary}
\begin{proof}
Apply Theorem \ref{closure} repeatedly and then
Corollary \ref{arithmetic}.
\end{proof}

For the obstruction below use the ordinary unitary affine
$F_4$ category at level four, realized by a completely
rational conformal net by \cite[Theorem I and Theorem 3.40]{Gui}.
Its object $a=X_5$ has highest weight $2\lambda_4$,
conformal weight one and $18<d_a<19$.
The total matched-channel multiplicities in $a^*X_x$
for $x=0,\ldots,15$ are
\[
 (1,1,1,1,1,3,2,2,2,2,2,3,3,3,3,4).
\]
Thus $D_a(x)>0$ for every simple probe, and $d_a$ is
nonintegral. Theorem \ref{permutation-main} applies.
All of its prime permutation gaugings also satisfy joint
detection by Theorem \ref{closure}. This separates the
joint condition from the sufficient one-probe inequality.


\section{Permutation amplification of the obstruction}
Let $\mathcal A$ be a completely rational unitary conformal net,
and let $\C=\operatorname{Rep}(\mathcal A)$ with its positive ribbon
structure. Write $d_x$ for dimensions, $\theta_x$ for twists,
$S$ for the normalized positive Hopf matrix, and $\Delta=\dim\C$.
For a simple boson $a$, put
\[
 D_a(x)=\sum_{b:\,\theta_b=\theta_x}N_{a^*,x}^b d_b.
\]
By balancing, this is the dimension of
the monodromy-one summand in $a^*x$.

\begin{theorem}\label{permutation-main}
Suppose $a\in\Irr(\C)$ satisfies $\theta_a=1$, $d_a>1$,
and
\[
 D_a(x)>0\quad\text{for every }x\in\Irr(\C).
                                                               \tag{P1}
\]
For every sufficiently large prime $p$, the cyclic permutation
orbifold category
\[
 \E_p=\operatorname{Rep}
       \bigl((\mathcal A^{\otimes p})^{\langle(1\cdots p)\rangle}\bigr)
\]
has a simple boson $A_p$ of dimension $d_a^p$ with
\[
               D_{A_p}(X)\ge d_{A_p}\quad(X\in\Irr(\E_p)).        \tag{P2}
\]
Equality holds exactly for the character lifts of the constant
neutral tuples $(u,\ldots,u)$ with $u$ invertible in $\C$.
If $d_a$ is nonintegral, then $d_{A_p}$ is nonintegral for
all but at most one prime $p$.
\end{theorem}

The $F_4$ object in Section \ref{sec:seeds} satisfies (P1). Thus the theorem gives infinitely many examples of nonintegral bosons with no contracting simple probe.

\subsection{Prime-cycle formulas and the diagonal lift}
The orbifold is completely rational by
\cite[Theorem 2.7]{Xu}. Its neutral simple objects are the orbits
of nonconstant $p$-tuples, of orbit size $p$, and the $p$ character
lifts of each constant tuple. Each nonzero degree $g^r$ has
simples $X=(y,g^r,k)$, with $y\in\Irr(\C)$ and
$k\in\mathbb Z/p$. These are the prime-cycle cases of
\cite[Section 2.12]{Xu} and
\cite[Proposition 6.7]{DRX}.

Choose the trivial-character diagonal lift $A_p=(a,1)$
in the notation of \cite[Lemma 2.12]{Xu}. Its underlying
neutral object is $(a,\ldots,a)$, so its dimension and
twist are $d_a^p$ and one. For odd $p$ the automorphism
$\tau_p$ in that lemma is the unit by
\cite[Proposition 3.4(2)]{Xu}; no hypothesis on the pointed
part of $\C$ is needed. Consequently
\[
 S^{\E_p}_{A_p,(y,g^r,k)}=S_{a,y}/p\qquad(r\ne0).                 \tag{P3}
\]
Indeed Lemma 2.12 gives this at $r=1$, and conjugating the
cycle by a tensor permutation gives every $r\ne0$ while
preserving the trivial-character diagonal lift. The same
permutation argument occurs in the proof of
\cite[Proposition 3.8(2)]{Xu}.

The ratios of the $A_p$ row to the positive vacuum row
are products of $r_a(z)=S_{a,z}/S_{0,z}$ on neutral
tuples, and are $r_a(z)$ on twisted columns by (P3).
Put $b=a^*$, so $r_b(z)=\overline{r_a(z)}$ and
$\theta_b=1$. The ratios for $A_p^*$ are consequently
products of $r_b(z)$ and the individual $r_b(z)$,
respectively, by duality of the normalized modular matrix.
Moreover
\[
 d_{(y,g^r,k)}=\Delta^{(p-1)/2}d_y,                              \tag{P4}
\]
using (P3) with $a=\one$ and
$S^{\E_p}_{\one,\one}=S_{0,0}^{\,p}/p$.

We will use the following prime-cycle matrix formulas from
\cite[Proposition 3.8 and Theorem 3.12]{Xu}.
Xu's inverse Hopf convention is translated by conjugating
the base and orbifold matrices simultaneously. The mixed
formula (P3) retains its form, and all entrywise absolute
squares below are unchanged; the translated $\Lambda$
matrices remain unitary.
A twisted row vanishes on nonconstant neutral columns; on
$(z,\ldots,z,l)$ its squared modulus is $|S_{y,z}|^2/p^2$.
On $(z,g^s,l)$ its squared modulus is
\[
          p^{-2}|\Lambda(sr^{-1}/p)_{y,z}|^2.                    \tag{P5}
\]
Here the quotient $sr^{-1}$ is read modulo $p$, and the
$\Lambda$ matrices are unitary. The difference between
Xu's matrices and these unitary matrices is a scalar
phase; hence it disappears in (P5).

\subsection{Uniform neutral-sector inequality}
Let $H$ be the finite cyclic group containing all twist ratios
in $\C$. For a simple $x$ set
\[
 w_x=\frac1{d_a}\sum_c N_{b,x}^c d_c[\theta_c/\theta_x]
                      \ \in\mathbb R_{\ge0}[H].                 \tag{P6}
\]
Its total mass is $d_x$, and its identity coefficient is
$D_a(x)/d_a>0$. If $x=u$ is invertible, $bu$ is simple;
the positive identity coefficient forces
$\theta_{bu}/\theta_u=1$. Since $d_{bu}=d_a$, this gives
$w_u=[1]$.

\begin{lemma}\label{coefficient}
For the finitely many noninvertible labels $x$, put
$c_x(n)=[1]w_x^n$, with $c_x(0)=1$. Then each $c_x(n)>0$,
each $\inf_{n\ge0}c_x(n)>0$, and $c_x(n)\to\infty$.
There is an integer $K$ such that
\[
 [1]\prod_x w_x^{n_x}>1
       \quad\text{if}\quad\sum_{x\text{ noninvertible}}n_x\ge K.
                                                               \tag{P7}
\]
For smaller sums these coefficients have a positive lower bound
$\epsilon$ independent of all invertible exponents.
\end{lemma}
\begin{proof}
Fix a noninvertible $x$ and let $H_x\le H$ be generated by
the support of the probability measure $\mu_x=w_x/d_x$.
Its positive identity mass implies that every nontrivial
character $\chi$ of $H_x$ satisfies
$|\widehat\mu_x(\chi)|<1$: equality in the triangle
inequality would make $\chi$ constant on the support,
and its identity value would make that constant one,
contrary to generation. Finite Fourier inversion gives
\[
  [1]\mu_x^n=\frac1{|H_x|}
            \sum_{\chi\in\widehat H_x}\widehat\mu_x(\chi)^n
            \longrightarrow\frac1{|H_x|}.
\]
Since $d_x>1$, $c_x(n)\to\infty$. Every $c_x(n)$ is
positive because the identity belongs to the support.
Its infimum is positive by its limit and finitely many
initial terms.

The identity coefficient in a product is at least the
product of the identity coefficients. If the sum of the
noninvertible exponents tends to infinity, some exponent
tends uniformly to infinity after division by the fixed
number of noninvertible labels. Its factor $c_x(n_x)$
tends uniformly to infinity, while all other factors
have fixed positive lower bounds. This proves (P7).
Below $K$ there are only finitely many such exponent tuples,
all giving positive coefficients. Invertible factors have
$w_u=[1]$ and do not change the coefficient.
\end{proof}

For a neutral probe $X$ above $(x_1,\ldots,x_p)$,
forgetting equivariance and using the neutral braiding gives
\[
 \frac{D_{A_p}(X)}{d_a^p}
   =\begin{cases}
      p[1]\prod_{i=1}^p w_{x_i},&\text{nonconstant tuple},\\
      [1]w_x^p,&\text{constant tuple }(x,\ldots,x).
     \end{cases}                                               \tag{P8}
\]
Choose $p\ge K$ and $p\epsilon>1$.
For a nonconstant tuple with at least $K$ noninvertible
entries, (P7) and (P8) give strict inequality in (P2).
With fewer entries, the lower bound is $p\epsilon>1$.
A constant noninvertible tuple has $p\ge K$ such entries
and also gives strict inequality. A constant invertible
tuple has ratio one for each character lift. This proves
the claimed neutral result, including its equality cases.

\subsection{Uniform twisted-sector inequality}
Verlinde and formulas (P3)--(P5) give, for $X=(y,g^r,k)$,
\[
 N_{A_p^*,X}^{X}
  =\frac1p\left((N_b^p)_{y,y}
       +\sum_{s=1}^{p-1}\sum_z r_b(z)
                    |\Lambda(sr^{-1}/p)_{y,z}|^2\right).
                                                               \tag{P9}
\]
The constant neutral columns contribute the first term,
since there are $p$ characters and their row ratio is
$r_b(z)^p$. Each nonzero degree contributes the indicated
second term after its $p$ characters are summed.
The numbers $r_b(z)$ have absolute value at most
$d_a$, being eigenvalues of the nonnegative fusion matrix
$N_b$. Each weighted row sum has absolute value at most
$d_a$, by unitarity of $\Lambda$. Taking real parts in (P9)
therefore yields
\[
 N_{A_p^*,X}^{X}\ \ge\
      \frac{(N_b^p)_{y,y}-(p-1)d_a}{p}.                         \tag{P10}
\]

We now show that (P1) itself supplies sufficiently many unit
summands in $b^n$, without a return hypothesis. If
$|r_b(z)|=d_a$, balancing gives
\[
 \frac{r_b(z)}{d_a}
     =\sum_c\frac{N_{b,z}^c d_c}{d_a d_z}
                       \frac{\theta_c}{\theta_z}.
\]
This is a convex combination of unit complex numbers.
Its absolute value is one only when all terms with
positive weight are the same phase. By (P1) at least
one of those phases equals one, so $r_b(z)=d_a$.
Thus every peripheral eigenvalue of $N_b$ equals $d_a$.
Verlinde diagonalization by the unitary matrix $S$ gives
\[
 \frac{m_n}{d_a^n}:=
 \frac{N_{b^n}^{\one}}{d_a^n}
   =\sum_z |S_{0,z}|^2\left(\frac{r_b(z)}{d_a}\right)^n
     \longrightarrow c_a:=\sum_{z:r_b(z)=d_a}|S_{0,z}|^2>0.
                                                               \tag{P11}
\]
The last inequality holds because $r_b(0)=d_a$ and
$S_{0,0}>0$. Since $b^n$ contains $m_n$ units,
$(N_b^n)_{y,y}\ge m_n$ for every simple $y$ of $\C$.
For all sufficiently large $p$, (P11) gives
$m_p>(p-1)d_a$, so (P10) is positive for every twisted
simple, without any dependence on $y,r,k$.

Because $A_p$ is a boson, its dual's positive self-channel
is counted in $D_{A_p}(X)$. It follows from (P4) that
\[
 D_{A_p}(X)\ge d_X\ge\Delta^{(p-1)/2}>d_a^p
\]
for all sufficiently large $p$; the last inequality uses
$\Delta>d_a^2$. This proves the strict twisted result.

Finally, if $d_a^p=M$ and $d_a^q=N$ were integers for
two distinct primes, B\'ezout's identity would give
$d_a=M^uN^v\in\mathbb Q$ for integers $u,v$.
As a dimension it is an algebraic integer, hence an
integer. This proves the last assertion and completes
Theorem \ref{permutation-main}.


\begin{thebibliography}{99}
\bibitem{KT} C. Kassel and V. Turaev, \emph{Braid Groups}, Springer, 2008.
\bibitem{CGPW} S. X. Cui, C. Galindo, J. Y. Plavnik and Z. Wang,
\emph{On gauging symmetry of modular categories}, Comm. Math. Phys.
348 (2016), 1043--1064.
\bibitem{GN} J. Green and D. Nikshych,
\emph{On the braid group representations coming from weakly
group-theoretical fusion categories}, arXiv:1911.02633v1 (2019).
\bibitem{GS} P. Godfard, \emph{Semisimplicity of conformal blocks},
arXiv:2507.06318v2 (2026).
\bibitem{P56} Tethys, \emph{Vacuum braid constituents and good reduction},
accepted P56, version 6, Journal of Modular Tensor Categories (2026).
\bibitem{Atlas} J. Berean-Dutcher, C. C\'ordova, J. Parra-Martinez
and R. Tuddenham, \emph{The Anyon Atlas}, arXiv:2610.25439 (2026).
\bibitem{Gui} B. Gui, \emph{Unbounded field operators in categorical
extensions of conformal nets}, arXiv:2001.03095v5 (2026).
\bibitem{Xu} F. Xu, \emph{Some computations in the cyclic permutations
of completely rational nets}, arXiv:math/0511662v1 (2005).
\bibitem{DRX} C. Dong, L. Ren and F. Xu,
\emph{$S$-matrix in orbifold theory}, arXiv:2009.02612v1 (2020).
\end{thebibliography}

\end{document}
